\section{How General the Result Is}\label{sec:generality}

\Cref{thm:main} is not about one system. Discharging the section variables of the
mechanisation, it reads in full:
\begin{lstlisting}
forall (Opub Opriv : Ty) (runtime runs : nat)
       (p_pub : Proc Empty Opub)
       (p_priv : Proc THandlerOutput Opriv)
       (p_sched : Proc Empty Nat),
  NI (public_equiv Empty) (public_equiv Opub)
     p_pub ->
  NI (private_equiv THandlerOutput)
     (private_equiv Opriv) p_priv ->
  NI (public_equiv Empty) (public_equiv Nat)
     p_sched ->
  NI in_equiv (out_equiv Opub Opriv)
     (model_sliced runtime runs
                   p_pub p_priv p_sched)
\end{lstlisting}
Everything the pool is built around stays fixed; everything it carries is left
open (\cref{tab:open}). The side conditions are the expected ones: the mechanism
cannot make a leaking process secure, so each slot is assumed non-interfering at its
own classification, and the theorem says that composing them through this mechanism
leaks nothing further. The handler length and slice size carry no side condition,
because the slice is \emph{defined} as $\mathit{runs}\cdot\mathit{runtime}$ and so
ends on a handler boundary by construction.

\begin{table}[t]
\centering
\caption{What \cref{thm:main} leaves open.}
\label{tab:open}
\footnotesize
\begin{tabularx}{\columnwidth}{@{}lLl@{}}
\toprule
parameter & ranges over & side condition \\
\midrule
\code{p_pub} & any public user process & $\NI$ at \m{public_equiv} \\
\code{p_priv} & any private user process & $\NI$ at \m{private_equiv} \\
\code{p_sched} & any scheduler & $\NI$ at \m{public_equiv} \\
\code{Opub}, \code{Opriv} & the user processes' output alphabets & none \\
\code{runtime} & handler length & none \\
\code{runs} & handler runs per slice & none \\
\bottomrule
\end{tabularx}
\end{table}

Two features of the design make this possible. The state update never reads the
\emph{value} a user slot emits, only whether it emitted (the scheduler-output test
inspects the user slots for \code{None} and nothing more), so no user behaviour reaches
the schedule. And the state obligation \m{fv_NI} never mentions the slot processes, so
the one hard obligation is unaffected by what fills them.

What stays fixed is the mechanism under study, the three handlers, and the shape of
the pool around it: six slots in a fixed order. The order is hard-wired in the
projections that read the pool's output, and varying the number of user slots would
change the state transition and with it the proof of \m{fv_NI}.

\paragraph{Levels} We have read levels as two. The mechanisation works over an
arbitrary bounded lattice, where a characterised equivalence must also be monotone
downwards (what a level relates stays related below it, and likewise for
unobservability), and $\NI$ quantifies over every level. Nothing in the proof depends
on the lattice having two points, because every equivalence in the models
distinguishes $\bot$ from all other levels and nothing more. The level case splits
in the proof are $l = \bot$ against $l \neq \bot$, and the second case covers every
other level of any lattice at once.
